%!TEX root = forallxyyc.tex
\part{真值函项逻辑}
\label{ch.TFL}
\addtocontents{toc}{\protect\mbox{}\protect\hrulefill\par}

\chapter{迈向符号化的第一步}

\section{基于形式的有效性}\label{s:ValidityInVirtueOfForm}
考虑这个论证：


\begin{earg}
		\item[] 外面正在下雨。
		\item[] 如果外面正在下雨，那么 Jenny 很痛苦。
		\item[\texttherefore] Jenny 很痛苦。
	\end{earg}


再看另一个论证：


\begin{earg}
		\item[] Jenny 是一名无政府工团主义者。
		\item[] 如果 Jenny 是一名无政府工团主义者，那么 Dipan 是 Tolstoy 的忠实读者。
		\item[\texttherefore] Dipan 是 Tolstoy 的忠实读者。
	\end{earg}


这两个论证都是有效的，并且我们可以直截了当地说，它们共享一个共同的结构。我们可以将这个结构表达如下：


\begin{earg}
		\item[] $A$
		\item[] 如果 $A$，那么 $C$
		\item[\texttherefore] $C$
	\end{earg}


这看起来是一个出色的论证\emph{结构}。确实，任何具有这种结构的论证无疑都是有效的。但这并非唯一的良好论证结构。考虑如下论证：


\begin{earg}
		\item[] Jenny 要么开心要么难过。
		\item[] Jenny 不开心。
		\item[\texttherefore] Jenny 很难过。
	\end{earg}


同样，这是一个有效论证。其结构大致如下：


\begin{earg}
		\item[] $A$ 或 $B$
		\item[] 非 $A$
		\item[\texttherefore] $B$
	\end{earg}


一个绝佳的结构！这里还有另一个例子：


\begin{earg}
		\item[] Jim 并非既努力学习又参演了许多戏剧。
		\item[] Jim 努力学习。
		\item[\texttherefore] Jim 没有参演许多戏剧。
	\end{earg}


这个有效论证的结构或许可以这样表示：


\begin{earg}
		\item[] 非($A$ 和 $B$)
		\item[] $A$
		\item[\texttherefore] 非 $B$
	\end{earg}


这些例子阐明了一个重要的思想，我们可以称之为\emph{基于形式的有效性}。我们刚才所考虑论证的有效性，与“Jenny 很痛苦”、“Dipan 是 Tolstoy 的忠实读者”或“Jim 参演了许多戏剧”这类英语表达式的意义几乎毫无关系。如果说它与意义有关，那也仅仅是与“和”、“或”、“非”以及“如果……那么……”这类短语的意义有关。

在 \crefrange{ch.TFL}{ch.NDTFL} 中，我们将发展一种形式语言，它允许我们以某种方式符号化许多论证，从而展示它们是基于其形式而有效的。这就是\emph{真值函项逻辑}的语言。\footnote{真值函项逻辑也常被称为（经典）“命题逻辑”或“语句逻辑”。}

\section{基于特殊原因的有效性}

在 \cref{s:FormalValidity}，我们首次介绍了形式有效性的概念，并将其与因所考虑反例的种类不同而产生的其他有效性进行了对比。值得重申的是，存在大量虽然有效但其有效性并非源于形式的论证。举个例子：


\begin{earg}
		\item[] Juanita 是一只雌狐。
		\item[\texttherefore] Juanita 是一只狐狸。
	\end{earg}


这个前提为真而结论为假是不可能的，因为“雌狐”（vixen）意思就是“雌性狐狸”（female fox）。因此，该论证是（概念上）有效的。这种有效性无法由该论证的形式来解释。为了看清这一点，我们可以给出一个形式相同但无效的论证，例如：


\begin{earg}
		\item[] Juanita 是一只雌狐。
		\item[\texttherefore] Juanita 是一座大教堂。
	\end{earg}


同样，考虑如下论证：


\begin{earg}
		\item[] 这座雕塑通体为绿色。
		\item[\texttherefore] 这座雕塑并非通体为红色。
	\end{earg}


同样，似乎不可能存在前提为真而结论为假的情况，因为没有任何事物可以同时通体为绿又通体为红。所以这个论证是有效的，但这里有一个形式相同却无效的论证：


\begin{earg}
		\item[] 这座雕塑通体为绿色。
		\item[\texttherefore] 这座雕塑并非通体闪亮。
	\end{earg}


这个论证是无效的，因为完全有可能同时通体为绿色且通体闪亮。（人们可以用优雅的闪亮绿色清漆来涂刷这座雕塑。）第一个论证的有效性很可能源于颜色（或颜色词）的相互作用方式；但无论是否如此，使其有效的并不仅仅是论证的\emph{形式}。

我们可以总结出以下重要启示：\emph{真值函项逻辑至多能帮助我们理解那些因其形式而有效的论证。}

\section{原子语句}

在 \cref{s:ValidityInVirtueOfForm} 中，我们通过用单个字母替换句子中的\emph{组成部分}来开始分离论证的形式。因此，在本节的第一个例子中，“外面正在下雨”是“如果外面正在下雨，那么 Jenny 很痛苦”的一个组成部分，我们将这个组成部分替换为“$A$”。

真值函项逻辑的人工语言将这一思路彻底贯彻。我们从一些 \emph{语句字母} 开始。它们将是构造更复杂语句的基本构件（它们是真值函项逻辑的“原子”语句）。我们将使用单个大写字母作为真值函项逻辑的语句字母。字母表虽然只有二十六个字母，但我们可能需要考虑的语句字母数量却没有限制。通过给字母添加下标，我们可以获得新的语句字母。因此，以下是五个不同的真值函项逻辑语句字母：
	$$A, P, P_1, P_2, A_{234}$$

我们将使用语句字母来表示，或者说\emph{符号化}，某些英语句子。为此，我们提供一个 \define{符号化约定}，例如：


\begin{ekey}
		\item[A] 外面正在下雨
		\item[C] Jenny 很痛苦
	\end{ekey}


这样做，并非要将这一符号化关系\emph{一劳永逸地}固定下来。我们只是说，在当下，我们将把真值函项逻辑的语句字母“$A$”视为对英语句子“外面正在下雨”的符号化，并把真值函项逻辑的语句字母“$C$”视为对英语句子“Jenny 很痛苦”的符号化。之后，当我们处理不同的句子或不同的论证时，完全可以提供一个新的符号化约定；比如说：


\begin{ekey}
		\item[A] Jenny 是一名无政府工团主义者
		\item[C] Dipan 是 Tolstoy 的忠实读者
	\end{ekey}


重要的是要理解：一个英语句子可能具有的任何结构，在用真值函项逻辑的语句字母将其符号化后都会丢失。从真值函项逻辑的角度看，一个语句字母就只是一个字母。它可以用来构建更复杂的语句，但它本身不能被分解。

\newglossaryentry{sentence letter}
{
name=语句字母,
description={用于表示真值函项逻辑中基本语句的字母}
}
\newglossaryentry{atomic sentence}
{
name=原子语句,
description={用于表示基本语句的表达式；在真值函项逻辑中是语句字母，在一阶逻辑中是后接名称的谓词符号}
}

\newglossaryentry{symbolization key}
{
name=符号化约定,
description={显示真值函项逻辑中哪些自然语言句子由哪些\glspl{sentence letter}表示的列表}
}

\chapter{联结词}
\label{s:TFLConnectives}

在上一章中，我们探讨了如何用真值函项逻辑的语句字母来符号化相当基本的自然语言句子。这引出了对“and”、“or”、“not”等英语表达式的处理需求。这些就是\emph{联结词}——它们可以用来从旧句子构成新句子。在真值函项逻辑中，我们将使用逻辑联结词从原子构件构建复杂语句。真值函项逻辑中有五个逻辑联结词。下表总结了它们，并将在本章节中加以解释。

\newglossaryentry{connective}
{
name=联结词,
description={真值函项逻辑中用于将\glspl{sentence letter}组合成更大语句的逻辑算子}
}


\begin{table}[h]
	\center
	\begin{tabular}{l l l}

	\textbf{符号}&\textbf{名称}&\textbf{粗略含义}\\
	\hline
	\enot&否定&“并非……”\\
	\eand&合取&“既……又……”\\
	\eor&析取&“或者……或者……”\\
	\eif&条件&“如果……那么……”\\
	\eiff&双条件&“……当且仅当……”\\

	\end{tabular}
	\end{table}

这些并非英语中令人感兴趣的全部联结词。其他还有例如“unless”、“neither \dots{} nor \dots”和“because”。我们将看到，前两者可以用我们将讨论的联结词来表达，而最后一个则不行。“because”与其他联结词不同，它不是\emph{真值函项的}。


\section{否定}

考虑如何符号化以下句子：


\begin{enumerate}
	\item\label{not1} 玛丽在巴塞罗那。
	\item\label{not2} 并非玛丽在巴塞罗那。
	\item\label{not3} 玛丽不在巴塞罗那。
	\end{enumerate}


为了符号化 \cref*{not1}，我们需要一个语句字母。可以提供如下符号化约定：


\begin{ekey}
		\item[B] 玛丽在巴塞罗那。
	\end{ekey}


由于 \cref*{not2} 显然与 \cref*{not1} 相关，我们不会想用一个完全不同的语句字母来符号化它。粗略地说，\cref*{not2}的意思就类似于“并非$B$”。为了符号化它，我们需要一个否定符号。我们将使用“\enot”。现在，我们可以用“$\enot B$”来符号化 \cref*{not2}。

\Cref*{not3} 也含有“不”这个词，并且它显然等价于 \cref*{not2}。因此，我们也可以用 “$\enot B$” 来符号化它。
\factoidbox{
	如果一个句子能用英语改述为“并非……”，那么它可以符号化为 $\enot\metav{A}$。
}
再多举几个例子会有帮助：


\begin{enumerate}
		\item\label{not4} 这个小部件可以被替换。
		\item\label{not5} 这个小部件是不可替代的。
		\item\label{not5b} 这个小部件不是不可替代的。
	\end{enumerate}


让我们使用下面的符号化约定：


\begin{ekey}
		\item[R] 这个小部件是可替代的。
	\end{ekey}


现在，\Cref*{not4} 可以符号化为 “$R$”。再看 \cref*{not5}：说这个小部件不可替代，意思就是“并非这个小部件是可替代的”。所以，即便 \cref*{not5} 不包含“不”这个词，我们也将它符号化为：“$\enot R$”。

\Cref*{not5b} 可以改述为“并非这个小部件是不可替代的”。而这又可以改述为“并非‘并非这个小部件是可替代的’”。因此，我们可以用真值函项逻辑语句 “$\enot\enot R$” 来符号化这个英语句子。

但是，在处理否定时需要特别小心。考虑：


\begin{enumerate}
		\item\label{not6} 简是快乐的。
		\item\label{not7} 简是不快乐的。
	\end{enumerate}


如果我们用真值函项逻辑语句 “$H$” 来符号化“简是快乐的”，那么我们可以把 \cref*{not6} 符号化为 “$H$”。然而，将 \cref*{not7} 符号化为 “$\enot{H}$” 将是一个错误。如果简是不快乐的，那么她确实不快乐；但 \cref*{not7} 的意思并不等同于“并非简是快乐的”。简可能既不快乐也非不快乐；她可能处于一种漠不关心的状态。因此，为了符号化 \cref*{not7}，我们就需要一个全新的真值函项逻辑语句字母。
\newglossaryentry{negation}
{
name=否定,
description={符号 \enot，用于表示在功能上类似于英语单词“not”的词语和短语}
}

\section{合取}
\label{s:ConnectiveConjunction}

考虑下面这些句子：


\begin{enumerate}
		\item\label{and1} 亚当是健壮的。
		\item\label{and2} 芭芭拉是健壮的。
		\item\label{and3} 亚当是健壮的，并且芭芭拉也是健壮的。
	\end{enumerate}


我们需要用不同的真值函项逻辑语句字母来符号化 \cref*{and1,and2}；可以这样设定：


\begin{ekey}
		\item[A] 亚当是健壮的。
		\item[B] 芭芭拉是健壮的。
	\end{ekey}


现在，\cref*{and1} 可以符号化为“$A$”，而 \cref*{and2} 可以符号化为“$B$”。\Cref*{and3} 大致说了“A 并且 B”。我们需要另一个符号来处理“并且”。我们将使用“$\eand$”。因此，我们将其符号化为“$(A\eand B)$”。这个联结词叫做 \define{合取}。我们也将“$A$”和“$B$”称为合取式“$(A \eand B)$”的两个 \define{合取支}。
\newglossaryentry{conjunction}
{
name=合取,
description={符号 \eand，用于表示在功能上类似于英语单词“and”的词语和短语；或是使用该符号构成的语句}
}

\newglossaryentry{conjunct}
{
name=合取支,
description={通过\gls{conjunction}与另一个语句连接在一起的语句}
}

注意，我们并不打算符号化 \cref*{and3} 中的“也”这个词。像“都”和“也”这类词的功能是提醒我们注意两个事物正在被合取。它们也许影响句子的强调意味，但在真值函项逻辑中，我们不会（也不能）符号化这类东西。

更多的例子将阐明这一点：


\begin{enumerate}
		\item\label{and4} 芭芭拉是健壮且有活力的。
		\item\label{and5} 芭芭拉和亚当都是健壮的。
		\item\label{and6} 虽然芭芭拉有活力，但她并不健壮。
	\item\label{and7} 亚当是健壮的，但芭芭拉比他更健壮。
	\end{enumerate}


\Cref*{and4} 显然是一个合取。这句话说了（关于芭芭拉的）两件事。在英语中，只提及一次芭芭拉是被允许的。我们\emph{可能}会忍不住想用类似于“”$B$” and energetic”这样的东西来符号化 \cref*{and4}。这会是一个错误。一旦我们将句子的一部分符号化为 “$B$”，任何进一步的结构就都丢失了，因为 “$B$” 是一个真值函项逻辑的语句字母。反过来说，“energetic”本身根本不是一个英语句子。我们的目标应该是像“”$B$” 并且芭芭拉是有活力的”这样的东西。所以我们需要向符号化约定中再添加一个语句字母。让 “$E$” 符号化“芭芭拉是有活力的”。现在，整个句子就可以符号化为 “$(B\eand E)$”。

\Cref*{and5} 对两个不同的主体说了同一件事。它声称芭芭拉和亚当都是健壮的，尽管在英语中“健壮的”这个词只用了一次。这个句子可以改述为“芭芭拉是健壮的，并且亚当是健壮的”。利用我们一直在使用的同一套符号化约定，我们可以用真值函项逻辑将其符号化为 “$(B\eand A)$”。

\Cref{and6} 略微复杂一些。 “although”这个词在句子的第一部分和第二部分之间建立了一种对比。尽管如此，这句话同时告诉我们芭芭拉精力充沛，而且她并不健壮。为了把每个合取支都变成一个语句字母，我们需要把“she”替换成“芭芭拉”。因此，我们可以将 \cref*{and6}  改写为：“\emph{既}芭芭拉精力充沛，\emph{又}芭芭拉不健壮”。第二个合取支包含一个否定，所以我们进一步改述：“\emph{既}芭芭拉精力充沛，\emph{又}\emph{并非}芭芭拉是健壮的”。现在我们可以用真值函项逻辑语句“$(E\eand\enot B)$”来符号化它。注意，在这个符号化过程中我们丢失了各种各样的细微之处。\cref*{and6} 和“既芭芭拉精力充沛，又并非芭芭拉是健壮的”在语气上有明显的不同。真值函项逻辑不能（也无法）保留这些细微差别。

\Cref{and7} 提出了类似的问题。
“but”这个词暗示了一种对比或差异，但这是真值函项逻辑无法处理的。我们只能把这个句子改述为“\emph{既}亚当是健壮的，\emph{又}芭芭拉比亚当更健壮”。（注意我们又一次把代词“him”替换成了“亚当”。）我们该如何处理第二个合取支呢？我们已经有语句字母“$A$”，它用来符号化“亚当是健壮的”，还有语句“$B$”用来符号化“芭芭拉是健壮的”；但这二者都不涉及他们的相对健壮程度。因此，为了符号化整个句子，我们需要一个新的语句字母。令真值函项逻辑语句“$R$”符号化英语句子“芭芭拉比亚当更健壮”。现在我们可以用“$(A \eand R)$”来符号化 \cref*{and7}。
\factoidbox{一个句子如果可以改述为“既……，又……”，或“……，但是……”，或“虽然……，……”，那么它就可以符号化为$(\metav{A}\eand\metav{B})$。}

我们之前提到，“but”所暗示的对比在真值函项逻辑中无法被捕捉，我们直接忽略它。一个不能简单忽略的现象是时间顺序。例如，考虑：

\begin{enumerate}
	\item\label{stoodfirst} Harry 站起来并反对该提议。
	\item\label{stoodsecond} Harry 反对该提议并站起来。
\end{enumerate}

如果 Harry 在反对\emph{之后}才站起来，那么 \cref*{stoodsecond} 为真而 \cref*{stoodfirst} 为假——这里“and”的使用是\emph{不对称的}。然而，真值函项逻辑的符号“\eand”总是对称的（或者像逻辑学家说的，满足“交换律”）。真值函项逻辑无法处理不对称的“and”。在本书所有的例子和练习中，我们都将假设“and”是对称的。\footnote{关于语言学中对称和不对称的合取，参见，例如，Robin Lakoff, “”If's, and's, and but's about conjunction'', 载于：C. J. Fillmore 和 D. T. Langendoen (编), \textit{Studies in Linguistic Semantics}, 霍尔特、莱因哈特与温斯顿出版社, 1971, 以及 Susan Schmerling, “”Asymmetric conjunction and rules of conversation'', 载于：P. Cole 和 J. L. Morgan (编), \textit{Speech Acts}, Brill, 1975, 第 211--31 页。}

你可能想知道为什么我们要在合取式周围加上括号。通过思考否定如何与合取交互，可以揭示其原因。考虑：

\begin{enumerate}
		\item\label{negcon1} 并非你既能得到汤又能得到沙拉。
		\item\label{negcon2} 你不会得到汤，但你会得到沙拉。
	\end{enumerate}

\Cref*{negcon1} 可以改述为“并非：你既会得到汤又会得到沙拉”。使用以下符号化约定：

\begin{ekey}
		\item[\ensuremath{S_1}] 你会得到汤。
		\item[\ensuremath{S_2}] 你会得到沙拉。
	\end{ekey}

我们会把“你既会得到汤又会得到沙拉”符号化为“$(S_1 \eand S_2)$”。那么，为了符号化 \Cref*{negcon1}，我们只需否定整个句子，即：“$\enot (S_1 \eand S_2)$”。

\Cref*{negcon2} 是一个合取式：你\emph{不会}得到汤，并且你\emph{会}得到沙拉。“你不会得到汤”被符号化为“$\enot S_1$”。因此，为了符号化 \Cref*{negcon2} 本身，我们给出“$(\enot S_1 \eand S_2)$”。

这些英语句子非常不同，它们的符号化也因此不同。在其中一个句子中，整个合取式被否定。在另一个中，只有一个合取支被否定。括号帮助我们厘清诸如否定的\emph{辖域}之类的东西。

\section{析取}

考虑以下句子：

\begin{enumerate}
		\item\label{or1}要么 Fatima 玩电子游戏，要么她看电影。
		\item\label{or2}要么 Fatima 要么 Omar 会玩电子游戏。
	\end{enumerate}

对于这些句子，我们可以使用以下符号化约定：

\begin{ekey}
		\item[F] Fatima 会玩电子游戏
		\item[O] Omar 会玩电子游戏
		\item[M] Fatima 会看电影
	\end{ekey}

然而，我们需要再次引入一个新符号。\Cref*{or1} 被符号化为“$(F \eor M)$”。这个联结词叫做 \define{析取（disjunction）}。我们也可以说，“$F$”和“$M$”是析取式“$(F \eor M)$”的 \define{析取支（disjuncts）}。

\newglossaryentry{disjunction}
{
name=析取,
description={联结词 \eor，用于表示功能类似于英语中“or”一词在相容意义上的词语和短语；或者指使用这个联结词构成的语句}
}

\newglossaryentry{disjunct}
{
name=析取支,
description={通过\gls{disjunction}联结到另一个语句的语句}
}

\Cref*{or2} 只是稍微复杂一点。它有两个主语，但原英语句子只给了一次动词。不过，我们可以将 \cref*{or2} 改述为：“要么 Fatima 会玩电子游戏，要么 Omar 会玩电子游戏”。现在我们显然可以再次将其符号化为“$(F \eor O)$”。
	\factoidbox{
		如果一个句子可以用英语改述为“要么……，要么……”，那么它可以符号化为 $(\metav{A}\eor\metav{B})$。
	}
有时在英语中，“or”一词的用法会排除两个析取支都为真的可能性。这叫做 \define{排斥或（exclusive or）}。当餐厅菜单上写着“主菜配汤或沙拉”时，显然就是指\emph{排斥或}：你可以要汤；你可以要沙拉；但如果你\emph{同时}想要汤\emph{和}沙拉，你就得额外付钱。

其他时候，“or”一词允许两个析取支可能都为真。上文中的 \cref{or2} 很可能就是这种情况。Fatima 可能独自玩电子游戏，Omar 可能独自玩电子游戏，或者他们可能一起玩。\Cref*{or2} 只不过是在说，\emph{至少}他们中有一个人会玩电子游戏。这是一种 \define{相容或（inclusive or）}。真值函项逻辑符号“\eor”总是符号化\emph{相容或}。

观察否定如何与析取相互作用也会有所帮助。考虑以下句子：

\begin{enumerate}
		\item \label{or3} 要么你不会得到汤，要么你不会得到沙拉。
		\item \label{or4} 你既不会得到汤也不会得到沙拉。
		\item \label{or.xor} 你会得到汤或沙拉，但不会两者都得。
\end{enumerate}

使用与之前相同的符号化约定，\cref*{or3} 可以这样改述：“\emph{要么}你不会得到汤，\emph{要么}你不会得到沙拉”。为了在真值函项逻辑中符号化这一点，我们需要析取和否定。“你不会得到汤”被符号化为“$\enot S_1$”。“你不会得到沙拉”被符号化为“$\enot S_2$”。因此 \cref*{or3} 本身被符号化为“$(\enot S_1 \eor \enot S_2)$”。

\Cref{or4} 也需要否定。它可以改述为：“\emph{并非：}你得到汤或你得到沙拉”。由于这否定了整个析取式，我们用“$\enot (S_1 \eor S_2)$”来符号化 \cref*{or4}。

\Cref{or.xor} 是一个\emph{排斥或}。我们可以把这个句子分成两部分。第一部分说你会得到其中之一。我们用“$(S_1 \eor S_2)$”来符号化这部分。第二部分说你不会两者都得。我们可以将其改述为：“并非你既得到汤又得到沙拉”。使用否定和合取，我们用“$\enot(S_1 \eand S_2)$”来符号化这部分。现在只需要把两部分结合起来。如上所述，“but”通常可以用“$\eand$”符号化。所以 \cref*{or.xor} 可以符号化为“$((S_1 \eor S_2) \eand \enot(S_1 \eand S_2))$”。

这最后一个例子表明了一些重要的东西。尽管真值函项逻辑符号“\eor”总是符号化\emph{相容或}，但我们可以在{真值函项逻辑}中符号化\emph{排斥或}。我们只需要额外使用几个其他符号。

\section{条件}
考虑这些句子：

\begin{enumerate}
		\item \label{if1} 如果 Jean 在巴黎，那么 Jean 在法国。
		\item \label{if2} Jean 在法国仅当 Jean 在巴黎。
\end{enumerate}

让我们使用以下符号化约定：

\begin{ekey}
		\item[P] Jean 在巴黎
		\item[F] Jean 在法国
\end{ekey}

\Cref*{if1} 大致是这种形式：“如果 $P$，那么 $F$”。我们将使用符号“\eif”来符号化这种“如果……，那么……”的结构。因此我们将 \cref*{if1} 符号化为“$(P\eif F)$”。这个联结词叫做 \define{条件（conditional）}。这里，“$P$”叫做条件式“$(P \eif F)$”的 \define{前件（antecedent）}，“$F$”叫做 \define{后件（consequent）}。

\newglossaryentry{conditional}
{
name=条件,
description={符号 \eif，用于表示功能类似于英语短语“如果……那么……”的词语和短语；或指使用这个符号构成的语句}
}

\newglossaryentry{antecedent}
{
name=前件,
description={\gls{conditional}左边的句子}
}

\newglossaryentry{consequent}
{
name=后件,
description={\gls{conditional}右边的句子}
}

\Cref{if2} 也是一个条件句。由于“if”这个词出现在句子的后半部分，可能会诱使人想用和 \cref*{if1} 同样的方式把它符号化。但这是错误的。你的地理知识告诉你 \cref*{if1} 显然是成立的：Jean 不可能在巴黎却不在法国。但 \cref*{if2} 就没这么显然了：如果 Jean 在迪耶普、里昂或者图卢兹，那么她虽然在法国却不在巴黎，这就会让 \cref*{if2} 为假。既然单靠地理知识就能决定 \cref*{if1} 的真假，而要判断 \cref*{if2} 的真假却需要（比方说）了解旅行计划，那么这两个句子的意思肯定不一样。

事实上，\cref*{if2} 可以改述为“如果 Jean 在法国，那么 Jean 在巴黎”。因此我们可以用 “$(F \eif P)$” 来符号化它。
	\factoidbox{
		如果一个句子可以用英语改述为“如果 A，那么 B”或“A 仅当 B”，那么它可以符号化为 $(\metav{A} \eif \metav{B})$。
	}
\noindent 事实上，条件句可以表示很多英语表达。考虑下面的句子：

\begin{enumerate}
		\item \label{ifnec1} Jean 要在巴黎，一个必要条件是她在法国。
		\item \label{ifnec2} Jean 在巴黎的必要条件是她身在法国。
		\item \label{ifsuf1} 要让 Jean 在法国，只要她在巴黎就足够了。
		\item \label{ifsuf2} Jean 在法国的充分条件是她身在巴黎。
\end{enumerate}

想一想就会发现，所有这四个句子的意思都和“如果 Jean 在巴黎，那么 Jean 在法国”一样。所以它们都可以符号化为 “$(P \eif F)$”。

要记住一点，联结词 “\eif” 只告诉我们：如果前件为真，那么后件也为真。它没有表达出两个事件之间（比方说）任何\emph{因果}联系。事实上，我们用 “$\eif$” 来符号化英语条件句时，会丢失大量信息。我们将在 \cref{s:IndicativeSubjunctive} 和~\cref{s:ParadoxesOfMaterialConditional} 中再讨论这个问题。

\section{双条件}
考虑下面这几个句子：

\begin{enumerate}
		\item \label{iff1} 只有当莱卡是哺乳动物时，它才是狗。
		\item \label{iff2} 只要莱卡是哺乳动物，它就是狗。
		\item \label{iff3} 莱卡是狗，当且仅当它是哺乳动物。
\end{enumerate}

我们将使用下面的符号化约定：

\begin{ekey}
		\item[D] 莱卡是狗
		\item[M] 莱卡是哺乳动物
\end{ekey}

根据前面讨论的理由，\cref*{iff1} 可以符号化为 “$(D \eif M)$”。

\Cref*{iff3} 表达的断言比 \cref*{iff1} 或 \cref*{iff2} 都要强。它可以改述为“如果莱卡是哺乳动物，那么莱卡是狗；并且，只有当莱卡是哺乳动物时，莱卡才是狗”。而这正是 \cref*{iff1,iff2} 的合取。因此我们可以把它符号化为 “$(D \eif M) \eand (M \eif D)$”。我们把它叫作\define{双条件（biconditional）}，因为它就是同时陈述了条件句的两个方向。

\newglossaryentry{biconditional}
{
name=双条件,
description={符号 \eiff，用于表示功能类似于英语短语“当且仅当”的词语和短语；或者由这个联结词构成的句子}
}

我们本可以用这种方式来处理所有双条件句。所以，就像我们不需要一个新的真值函项逻辑符号来处理\emph{排斥或}一样，我们其实也并非真的需要一个新符号来处理双条件句。只是因为双条件句出现得太频繁了，我们才决定使用符号 “\eiff”。这样，我们就可以用真值函项逻辑语句 “$(D \eiff M)$” 来符号化 \cref*{iff3}。

“当且仅当”这个表达在哲学、数学和逻辑学中尤其常见。为简洁起见，其英文原文“if and only if”常被缩写为“iff”。我们会遵循这个习惯。所以，一个“f”的“if”是英语的条件句。但两个“f”的“iff”是英语的双条件句。明确这一点后，我们就可以说：
	\factoidbox{
		如果一个句子可以用英语改述为“A \ifeff{} B”，即“A 当且仅当 B”，那么它可以符号化为 $(\metav{A} \eiff \metav{B})$。
	}
有一点要小心。日常说英语的人经常用“如果……，那么……”来表达他们真正想说的、更像是“……当且仅当……”的意思。也许在你小时候，你的父母对你说过：“如果你不吃掉青菜，你就没有甜点吃”。假设你吃掉了青菜，但你的父母却拒绝给你甜点，理由是：他们只承诺了那个\emph{条件句}（大致是“如果你吃了甜点，那你就肯定吃了青菜”），而不是双条件句（大致是“你吃到甜点，当且仅当你吃了青菜”）。那么，一场脾气发作恐怕就难免了。所以，在解读别人的话时要意识到这一点；但在你自己写作时，一定要确保只在你想表达\ifeff{}的双向意思时才使用双条件句。

\section{除非}
目前，我们已经介绍了真值函项逻辑的所有联结词。我们可以将它们结合起来，以符号化多种句子。一个特别困难的情形是使用英语联结词“unless”的时候：

\begin{enumerate}
\item\label{unless1} 除非你穿夹克，否则你会感冒。
\item\label{unless2} 你会感冒，除非你穿夹克。
\end{enumerate}

这两个句子显然是等价的。为了符号化它们，我们使用以下符号化约定：


\begin{ekey}
		\item[J] 你会穿夹克
		\item[D] 你会感冒
	\end{ekey}

这两个句子都意味着：如果你不穿夹克，那么你就会感冒。据此，我们可以将它们符号化为 “$(\enot J \eif D)$”。

同样，这两个句子也都意味着：如果你没有感冒，那么你一定穿了夹克。因此，我们可以将它们符号化为 “$(\enot D \eif J)$”。

同样，这两个句子也都意味着：你要么会穿夹克，要么就会感冒。基于此，我们可以将它们符号化为 “$(J \eor D)$”。

这三种符号化都是正确的。事实上，在\cref{s:SemanticConcepts}中我们将会看到，在真值函项逻辑中，这三种符号化是等价的。
% TODO: it might be useful to reference exercise 11.F.3 explicitly
% here, since the point is not discussed in the main text
	\factoidbox{
		如果一个句子可以改述为“除非 $A$，否则 $B$”，那么它可以符号化为 “$(\enot\metav{A}\eif\metav{B})$” 或 “$(\metav{A}\eor\metav{B})$”。
	}
不过，这里也有一点复杂之处。“除非”可以被符号化为一个条件句；但正如我们之前所说，人们经常在单独使用条件句时，想要表达的是双条件句。同样，“除非”也可以被符号化为一个析取句；但析取有两种（排斥析取和相容析取）。所以，日常说英语的人常常用“除非”来表示更接近双条件句或排斥析取的意思，这或许不会令你感到意外。假设有人说：“I will go running unless it rains”。他们的意思很可能是“当且仅当不下雨时，我才会去跑步”（即双条件句），或者“我要么去跑步，要么天会下雨，但两者不会同时成立”（即排斥析取）。再次提醒：解释他人的话语时需注意这一点，但你自己写作时应力求精确。

\practiceproblems
\solutions
\problempart 使用给定的符号化约定，将下列每个英语句子符号化为真值函项逻辑语句。\label{pr.monkeysuits}


\begin{ekey}
		\item[M] 那些生物是穿着西装的人
		\item[C] 那些生物是黑猩猩
		\item[G] 那些生物是大猩猩
	\end{ekey}

\begin{compactlist}
\item 那些生物不是穿着西装的人。
\item 那些生物是穿着西装的人，或者他们不是。
\item 那些生物要么是大猩猩，要么是黑猩猩。
\item 那些生物既不是大猩猩也不是黑猩猩。
\item 如果那些生物是黑猩猩，那么它们既不是大猩猩也不是穿着西装的人。
\item 除非那些生物是穿着西装的人，否则它们要么是黑猩猩，要么是大猩猩。
\end{compactlist}

\problempart 使用给定的符号化约定，将下列每个英语句子符号化为真值函项逻辑语句。


\begin{ekey}
\item[A] Ace 先生被谋杀了
\item[B] 管家干的
\item[C] 厨师干的
\item[D] 公爵夫人在撒谎
\item[E] Edge 先生被谋杀了
\item[F] 凶器是一个煎锅
\end{ekey}

\begin{compactlist}
\item 要么 Ace 先生，要么 Edge 先生被谋杀了。
\item 如果 Ace 先生被谋杀，那么是厨师干的。
\item 如果 Edge 先生被谋杀，那么不是厨师干的。
\item 要么是管家干的，要么公爵夫人在撒谎。
\item 只有在公爵夫人撒谎的情况下，才是厨师干的。
\item 如果凶器是一个煎锅，那么罪犯一定是厨师。
\item 如果凶器不是一个煎锅，那么罪犯要么是厨师，要么是管家。
\item Ace 先生被谋杀，当且仅当 Edge 先生没有被谋杀。
\item 公爵夫人在撒谎，除非被杀的是 Edge 先生。
\item 如果 Ace 先生被谋杀，他是被煎锅干掉的。
\item 既然厨师是凶手，那么管家就不是。
\item 公爵夫人当然在撒谎！
\end{compactlist}


\solutions

\problempart 使用给定的符号化约定，将下列每个英语句子符号化为真值函项逻辑语句。\label{pr.avacareer}


\begin{ekey}
		\item[\ensuremath{E_1}] Ava 是一名电工
		\item[\ensuremath{E_2}] Harrison 是一名电工
		\item[\ensuremath{F_1}] Ava 是一名消防员
		\item[\ensuremath{F_2}] Harrison 是一名消防员
		\item[\ensuremath{S_1}] Ava 对她的职业很满意
		\item[\ensuremath{S_2}] Harrison 对他的职业很满意
	\end{ekey}

\begin{compactlist}
\item Ava 和 Harrison 都是电工。
\item 如果 Ava 是消防员，那么她对自己的职业感到满意。
\item Ava 是消防员，除非她是电工。
\item Harrison 是一个不满意的电工。
\item Ava 和 Harrison 都不是电工。
\item Ava 和 Harrison 都是电工，但他们都对其不满意。
\item Harrison 感到满意仅当他是消防员。
\item 如果 Ava 不是电工，那么 Harrison 也不是；但如果她是，那么他也是。
\item Ava 对自己的职业感到满意当且仅当 Harrison 对自己的职业不满意。
\item 如果 Harrison 既是电工又是消防员，那么他必定对自己的工作感到满意。
\item Harrison 不可能既是电工又是消防员。
\item Harrison 和 Ava 都是消防员当且仅当他们俩都不是电工。
\end{compactlist}

\problempart
使用给出的符号化约定，将下列每个自然语言句子用真值函项逻辑符号化。
\label{pr.jazzinstruments}

\begin{ekey}
\item[\ensuremath{J_1}] John Coltrane 演奏了次中音萨克斯
\item[\ensuremath{J_2}] John Coltrane 演奏了高音萨克斯
\item[\ensuremath{J_3}] John Coltrane 演奏了大号
\item[\ensuremath{M_1}] Miles Davis 演奏了小号
\item[\ensuremath{M_2}] Miles Davis 演奏了大号
\end{ekey}

\begin{compactlist}
\item John Coltrane 演奏了次中音萨克斯和高音萨克斯。%{\color{red} $J_1 \eand J_2$} \vspace{1ex}
\item Miles Davis 和 John Coltrane 都没有演奏大号。%{\color{red} $\enot(M_2 \eor J_3)$ or $\enot M_2 \eand \enot J_3$} \vspace{1ex}
\item John Coltrane 没有同时演奏次中音萨克斯和大号。%{\color{red} $\enot(J_1 \eand J_3)$ or $\enot J_1 \eor \enotJ_3$} \vspace{1ex}
\item John Coltrane 没有演奏次中音萨克斯，除非他也演奏了高音萨克斯。%{\color{red} $\enot J_1 \eor J_2$} \vspace{1ex}
\item John Coltrane 没有演奏大号，但 Miles Davis 演奏了。%{\color{red} $\enotJ_3 \eand M_2$} \vspace{1ex}
\item Miles Davis 演奏小号仅当他也演奏大号。%{\color{red} $M_1 \eiff M_2$} \vspace{1ex}
\item 如果 Miles Davis 演奏了小号，那么 John Coltrane 至少演奏了这三种乐器中的一种：次中音萨克斯、高音萨克斯或大号。%{\color{red} $M_1 \eif (J_1 \eor (J_2 \eor J_3))&} \vspace{1ex}
\item 如果 John Coltrane 演奏了大号，那么 Miles Davis 既没有演奏小号也没有演奏大号。%{\color{red} $J_3 \eif \enot(M_1 \eor M_2)$ or $J_3 \eif (\enot M_1 \eand \enot M_2)$  } \vspace{1ex}
\item Miles Davis 和 John Coltrane 都演奏了大号当且仅当 Coltrane 没有演奏次中音萨克斯且 Miles Davis 没有演奏小号。%{\color{red} $(J_3 \eand M_2) \eiff \enotJ_1 & \enot M_1)$ or $(J_3 \eand M_2) \eiff \enot (J_1 \eor M_1)$} \vspace{1ex}
\end{compactlist}

\solutions
\problempart
\label{pr.spies}
给出一套符号化约定，并用真值函项逻辑将下列自然语言句子符号化。

\begin{compactlist}
\item Alice 和 Bob 都是间谍。
\item 如果 Alice 或 Bob 是间谍，那么密码被破解了。
\item 如果 Alice 和 Bob 都不是间谍，那么密码未被破解。
\item 德国大使馆将陷入混乱，除非有人破解了密码。
\item 要么密码被破解了，要么没有，但无论如何德国大使馆都会陷入混乱。
\item Alice 或 Bob 是间谍，但并非两人都是。
\end{compactlist}

\solutions
\problempart
给出一套符号化约定，并用真值函项逻辑将下列自然语言句子符号化。

\begin{compactlist}
\item 如果荣耀王国能找到食物，那么 Rafiki 会谈论压扁的香蕉。
\item Rafiki 会谈论压扁的香蕉，除非 Simba 还活着。
\item 要么 Rafiki 会谈论压扁的香蕉，要么不会，但无论如何荣耀王国都能找到食物。
\item Scar 将继续当国王当且仅当荣耀王国能找到食物。
\item 如果 Simba 还活着，那么 Scar 将不会继续当国王。
\end{compactlist}

\problempart
对于每个论证，写出一套符号化约定，并用真值函项逻辑将论证中的所有句子符号化。

\begin{compactlist}
\item 如果 Dorothy 在早上弹钢琴，那么 Roger 醒来时会脾气暴躁。Dorothy 在早上弹钢琴，除非她分心。所以，如果 Roger 醒来时没有脾气暴躁，那么 Dorothy 一定是分心了。
\item 星期二要么下雨要么下雪。如果下雨，Neville 会难过。如果下雪，Neville 会觉得冷。因此，星期二 Neville 要么难过要么觉得冷。
\item 如果 Zoog 记得做家务，那么房间是干净的但不整洁。如果他忘记了，那么房间是整洁的但不干净。因此，房间要么整洁要么干净，但不会同时既整洁又干净。
\end{compactlist}

\problempart
对于每个论证，写出一套符号化约定，并尽可能用真值函项逻辑将论证符号化。段落中\emph{斜体字部分}是为论证提供语境，不需要符号化。

\begin{compactlist}
\item 很快就会下雨。我知道是因为我的腿在疼，而如果要下雨的话，我的腿就会疼。

\item  \emph{蜘蛛侠试图搞清楚坏蛋的计划。}如果章鱼博士得到了铀，他就会勒索这座城市。我对此确信无疑，因为如果章鱼博士得到了铀，他就能造出脏弹；而如果他造出了脏弹，他就会勒索这座城市。

\item \emph{一位西方人试图预测中国政府的政策。}如果中国政府不能解决北京的水资源短缺问题，他们将不得不迁都。他们不想迁都。因此他们必须解决水资源短缺问题。但解决水资源短缺的唯一方法是从长江调走几乎全部的水资源北上。因此中国政府会继续进行南水北调工程。

\end{compactlist}

\problempart
我们曾使用“$\eor$”、“$\eand$”和“$\enot$”来符号化一个\emph{排斥或}。你能否只用两个联结词来符号化一个\emph{排斥或}？有没有可能只用唯一一个联结词来符号化一个\emph{排斥或}呢？

\chapter{真值函项逻辑的语句}\label{s:TFLSentences}
“要么苹果是红色的，要么浆果是蓝色的”是一个英语句子，而“$(A\eor B)$”是一个真值函项逻辑语句。虽然我们在遇到英语句子时能够认出它们，但我们并没有对“英语句子”下一个形式定义。然而，在这一章中，我们将精确定义什么才算是一个真值函项逻辑语句。这就是像真值函项逻辑这样的形式语言比像英语这样的自然语言更精确的一个方面。

\section{表达式}

我们已经看到，真值函项逻辑中有三种符号：

\begin{center}
\begin{tabular}{l l}
语句字母 & $A,B,C,\ldots,Z,\ldots$ \\
（可根据需要加下标） & $A_1, B_1,Z_1,A_2,A_{25},J_{375},\ldots$ \\
\\
联结词 & $\enot,\eand,\eor,\eif,\eiff$ \\
\\
括号 &( , ) \\
\end{tabular}
\end{center}

定义一个\define{真值函项逻辑表达式}为任意一串真值函项逻辑的符号。
所以：以任何顺序写下真值函项逻辑的任意符号序列，你就得到了一个真值函项逻辑表达式。\footnote{在\cref{sec:contradiction}中，我们将引入符号“$\ered$”，它也将被视为一个原子语句。}

\section{语句}\label{s:Sentences}
根据我们刚刚所说的，“$(A \eand B)$”是一个真值函项逻辑表达式，“$\lnot)(\eor()\eand(\enot\enot())((B$”同样也是。但前者是一个\emph{语句}，而后者只是\emph{乱码}。我们需要一些规则来告诉我们哪些真值函项逻辑表达式是语句。

显然，单个的语句字母如“$A$”和“$G_{13}$”应该算作语句。（我们也称它们为\define{原子}语句。）我们可以使用各种联结词，由这些原子语句构成更多的语句。利用否定，我们可以得到“$\enot A$”和“$\enot G_{13}$”。利用合取，我们可以得到“$(A \eand G_{13})$”、“$(G_{13} \eand A)$”、“$(A \eand A)$”和“$(G_{13} \eand G_{13})$”。我们也可以通过重复运用否定得到像“$\enot \enot A$”这样的语句，或者将否定与合取结合得到像“$\enot(A \eand G_{13})$”和“$\enot(G_{13} \eand \enot G_{13})$”这样的语句。可能的组合是无穷无尽的，即便我们只从这两个语句字母开始；并且语句字母本身就有无穷多个。因此，试图逐个列举出所有语句是毫无意义的。

相反，我们将描述语句可以被\emph{构造}出来的过程。考虑否定：对于任意给定的真值函项逻辑语句\metav{A}，$\enot\metav{A}$ 都是真值函项逻辑的一个语句。（为什么使用这种特殊的字体？我们将在\cref{s:Metavariables}中回到这个问题。）

我们可以对其余每一个联结词给出类似的说明。例如，如果 \metav{A} 和 \metav{B} 是真值函项逻辑的语句，那么$(\metav{A}\eand\metav{B})$就是一个真值函项逻辑的语句。为所有联结词提供这样的条款后，我们得到了下面关于\define{真值函项逻辑语句}的形式定义：
	\factoidbox{\label{TFLsentences}


\begin{numberlist}[labelindent=0pt,leftmargin=*]
		\item 每一个语句字母都是一个语句。
		\item 如果 \metav{A} 是一个语句，那么 $\enot\metav{A}$ 就是一个语句。
		\item 如果 \metav{A} 和 \metav{B} 是语句，那么$(\metav{A}\eand\metav{B})$ 就是一个语句。
		\item 如果 \metav{A} 和 \metav{B} 是语句，那么$(\metav{A}\eor\metav{B})$ 就是一个语句。
		\item 如果 \metav{A} 和 \metav{B} 是语句，那么$(\metav{A}\eif\metav{B})$ 就是一个语句。
		\item 如果 \metav{A} 和 \metav{B} 是语句，那么$(\metav{A}\eiff\metav{B})$ 就是一个语句。
		\item 除此之外，没有别的东西是语句。
	\end{numberlist}

}
\newglossaryentry{sentence of TFL}
{
name=真值函项逻辑语句,
text = 真值函项逻辑语句,
plural = 真值函项逻辑语句,
description={一个依据\ifHTMLtarget \cref{s:Sentences}\else 第~\pageref{TFLsentences}页\fi 给出的归纳规则构建起来的真值函项逻辑符号串}
}

这类定义被称为\define{归纳的}。归纳定义首先指明一些可明确的基本元素，然后提供一些方法，通过组合先前已确立的元素，来生成无限多的新元素。为了让你更好地理解什么是归纳定义，我们可以为\emph{我的祖先}这个概念本身给出一个归纳定义。我们先给出一个基始条款：


\begin{itemize}
		\item 我的父母是我的祖先。
	\end{itemize}


然后我们再给出进一步的条款，例如：


\begin{itemize}
		\item 如果 $x$ 是我的祖先，那么 $x$ 的父母也是我的祖先。
	\end{itemize}


最后，我们规定，\emph{只有}基始条款和归纳条款所说的才算是我的祖先。


\begin{itemize}
		\item 其他任何东西都不是我的祖先。
	\end{itemize}


使用这个定义，我们可以很容易地判断某人是否是我的祖先：只要检查她是否是我父母的……父母的父母就行了。对于我们给出的真值函项逻辑语句的归纳定义，情况也是一样。正如这个归纳定义允许复杂语句从较简单的部分构建起来一样，这个定义也允许我们把语句分解成更简单的部分。一旦我们分解到语句字母，分解过程就完成了。

让我们来看几个例子。

假设我们想知道 “$\enot \enot \enot D$” 是不是一个真值函项逻辑语句。根据定义的第二条款，如果 “$\enot \enot D$” 是一个语句，那么 “$\enot \enot \enot D$” 就是一个语句。所以现在我们需要问，“$\enot \enot D$” 是不是一个语句。再次查看定义的第二条款，如果 “$\enot D$” 是一个语句，那么 “$\enot \enot D$” 就是一个语句。同样，如果 “$D$” 是一个语句，那么 “$\enot D$” 就是一个语句。现在，“$D$” 是一个真值函项逻辑的语句字母，所以我们根据定义的第一条款知道“$D$”是一个语句。因此，对于像 “$\enot \enot \enot D$” 这样的复合语句，我们必须反复应用定义。最终我们会追溯到构建该语句所用的那些语句字母。

接下来，考虑例子 “$\enot (P \eand \enot (\enot Q \eor R))$”。根据定义的第二条款，如果“$(P \eand \enot (\enot Q \eor R))$”是一个语句，那么它就是语句；而如果\emph{既}有“$P$”\emph{又有}“$\enot (\enot Q \eor R)$”是语句，那么“$(P \eand \enot (\enot Q \eor R))$”就是一个语句。前者是一个语句字母，而后者若想成为语句，需要“$(\enot Q \eor R)$”是语句。它确实是。根据定义的第四条款，这又需要“$\enot Q$”和“$R$”都是语句，而它们的确都是！

归根结底，每个语句都是由语句字母规规矩矩地构造出来的。当我们处理\emph{除语句字母之外的语句}时，我们可以看到，在构造该语句时，必定有某个句子联结词是\emph{最后}被引入的。我们称这个联结词为该语句的\define{主逻辑算子}。就“$\enot\enot\enot D$”而言，主逻辑算子是最左边的那个“$\enot$”符号。就“$(P \eand \enot (\enot Q \eor R))$”而言，主逻辑算子是“$\eand$”。就“$((\enot E \eor F) \eif \enot\enot G)$”而言，主逻辑算子是“$\eif$”。

一般来说，你可以使用以下方法来找到一个语句的主逻辑算子：


\begin{itemize}
	\item 如果语句的第一个符号是“$\enot$”，那么它就是主逻辑算子。
	\item 否则，开始数括号。每遇到一个开括号“$($”就加 $1$；每遇到一个闭括号“$)$”就减 $1$。当你的计数恰好为 $1$ 时，你遇到的第一个算子（\emph{除了}“$\enot$”以外）就是主逻辑算子。
\end{itemize}

（注意：如果你使用这个方法，请确保包含语句中的\emph{所有}括号，而不是像\cref{TFLconventions}中约定的那样省略一些！）

真值函项逻辑语句的归纳结构，在我们考虑一个特定语句在什么情况下为真或为假时，将是重要的。语句“$\enot \enot \enot D$”为真，当且仅当语句“$\enot \enot D$”为假，如此顺着语句的结构继续分析，直到我们抵达其原子组成部分。我们将在\cref{ch.TruthTables}中回到这一点。

利用真值函项逻辑中语句的归纳结构，我们也可以给出否定\emph{辖域}的形式定义（见于\cref{s:ConnectiveConjunction}）。一个“$\enot$”的辖域就是该“$\enot$”作为主逻辑算子的子句。考虑如下语句：
$$(P \eand (\enot (R \eand B) \eiff Q))$$
该语句由将“$P$”与“$ (\enot (R \eand B) \eiff Q)$”进行合取而构成。后面这个语句则是通过在“$\enot (R \eand B)$”与“$Q$”之间放置双条件词构成的。前者——原语句的一个子句——是一个以“$\enot$”为主逻辑算子的语句。因此，该否定的辖域就是“$\enot(R \eand B)$”。更一般地说：
	\factoidbox{联结词（在一个语句中）的 \define{辖域（scope）} 就是该联结词作为主逻辑算子的子句。}

\section{括号约定}
\label{TFLconventions}
严格来说，“$(Q \eand R)$”中的括号是该语句不可或缺的一部分。部分原因在于，我们可能会将“$(Q \eand R)$”用作一个更复杂语句中的子句。例如，我们可能想要否定“$(Q \eand R)$”，从而得到“$\enot(Q \eand R)$”。如果我们只有不带括号的“$Q \eand R$”并在其前添加否定符号，就会得到“$\enot Q \eand R$”。我们很自然地会将其理解为与“$(\enot Q \eand R)$”意思相同，但正如我们在\cref{s:ConnectiveConjunction}中所见，这与“$\enot(Q\eand R)$”截然不同。

因此，严格地说，“$Q \eand R$”\emph{不是}一个语句。它仅是一个\emph{表达式}。

然而，在使用真值函项逻辑时，若我们偶尔不那么严格，可以使我们处理起来更方便。为此，我们采用以下一些便利约定。

首先，我们允许省略语句的\emph{最外层}括号。因此，我们允许用“$Q \eand R$”来代替语句“$(Q \eand R)$”。但必须记住，当我们要将该语句嵌入一个更复杂的语句时，需要把括号加回去！

其次，阅读带有大量嵌套括号的长句可能令人困扰。为使观感更轻松，我们允许使用方括号“[”和“]”来代替圆括号。例如，“$(P\eor Q)$”和“$[P\eor Q]$”在逻辑上没有区别。

结合这两条约定，我们可以将冗赘的语句
$$(((H \eif I) \eor (I \eif H)) \eand (J \eor K))$$
更清晰地改写如下：
$$\bigl[(H \eif I) \eor (I \eif H)\bigr] \eand (J \eor K)$$
现在，每个联结词的辖域就容易辨别得多了。

\practiceproblems

\solutions
\problempart
\label{pr.wiffTFL}
对以下各项：(a) 严格来说，它是真值函项逻辑的语句吗？(b) 按照我们放宽的括号约定，它是真值函项逻辑的语句吗？


\begin{compactlist}
\item $(A)$
\item $J_{374} \eor \enot J_{374}$
\item $\enot \enot \enot \enot F$
\item $\enot \eand S$
\item $(G \eand \enot G)$
\item $(A \eif (A \eand \enot F)) \eor (D \eiff E)$
\item $[(Z \eiff S) \eif W] \eand [J \eor X]$
\item $(F \eiff \enot D \eif J) \eor (C \eand D)$
\end{compactlist}

\problempart
是否存在不包含任何语句字母的真值函项逻辑语句？解释你的答案。\\

\problempart
该语句中每个联结词的辖域是什么？
$$\bigl[(H \eif I) \eor (I \eif H)\bigr] \eand (J \eor K)$$

\chapter{歧义性}\label{s:AmbiguityTFL}

在英语中，句子可能是 \define{歧义的（ambiguous）}，即它们可以具有多种含义。歧义有很多来源。一种是 \emph{词汇歧义：} 句子中包含具有多种含义的词语。例如，“bank”既可指河岸，也可指金融机构。因此，当我说“I went to the bank”时，可能意指我在河边散步，也可能意指我去存支票。根据不同的情境，“bank”一词的意图不同，因此该句子在不同的语境中说出口时，表达了不同的含义。

另一种歧义是 \emph{结构歧义}。当一个句子可以被不同方式解读时，便会产生这种歧义，且不同的解读对应不同的含义。诺姆·乔姆斯基的一个著名例子如下：

\begin{enumerate}
	\item[] Flying planes can be dangerous.
\end{enumerate}

一种解读是将“flying”用作修饰“planes”的形容词。在这种意义上，被断定为危险的是正在飞行中的飞机。在另一种解读中，“flying”是动名词：被断定为危险的是驾驶飞机的行为。在第一种情形下，你可能用这个句子来警告一个即将放飞热气球的人。在第二种情形下，你可能用它来劝某人不要当飞行员。

当句子被说出时，通常只有一种含义是说话者意图表达的。一个句子在具体说出时意图表达哪种可能含义，取决于语境，有时也取决于说话方式（例如，句子的哪些部分被重读）。通常，一种解读被意图的可能性远大于其他，以至于甚至很难“看出”那非意图的解读。这通常是笑话之所以好笑的原因，正如格劳乔·马克思的这个例子：

\begin{compactlist}
	\item[] 一天早晨，我穿着睡衣射杀了一头大象。
	\item[] 至于它怎么跑到我睡衣里的，我就不知道了。
\end{compactlist}

歧义性与模糊性相关，但并不相同。一个形容词，如“富有”或“高”，当无法总是判定它是否适用时，便是 \define{模糊的（vague）}。例如，一个身高 6~英尺 4~英寸（1.9~米）的人显然很高，但同样高度的建筑物却很矮小。在这里，语境对于判定哪些是明确适用、哪些是明确不适用的例子起着作用（“对人而言高”、“对篮球运动员而言高”、“对建筑物而言高”）。然而，即使语境明确，仍会存在处于中间区域的例子。

在真值函项逻辑中，我们通常力求避免歧义。我们会设法以不使用歧义词的方式给出符号化约定，或者，若一个词有多个含义，则消除其歧义。因此，例如，你的符号化约定需要为“丽贝卡去了（金融）银行”和“丽贝卡去了（河）岸”分配两个不同的语句字母。模糊性则更难避免。由于我们已经规定，在任何情况（以及后续，在任何赋值）下，每个基本语句（或语句字母）都必须为真或为假，不能处于中间状态，因此我们无法在真值函项逻辑中容纳边界情形。

真值函项逻辑语句的一个重要特征是它们\emph{不可能}是结构歧义的。每个真值函项逻辑语句都只能以一种，且仅此一种方式解读。真值函项逻辑的这一特性也是一个优点。若一个英语句子有歧义，真值函项逻辑能帮助我们澄明其不同的含义。尽管我们在日常对话中相当擅长处理歧义，但有时避免歧义至关重要。此时逻辑便能发挥有效作用：它帮助哲学家清晰表达思想，帮助数学家严谨陈述定理，帮助软件工程师明确地指定循环条件、数据库查询或验证标准。

在法律中，无歧义地陈述事项也至关重要。在此，毫不夸张地说，歧义可能关乎生死。以下是一个著名案例，其死刑判决取决于对法律中一处歧义的解释。罗杰·凯斯门特（1864--1916）是一位英国外交官，他因揭露刚果和秘鲁的侵犯人权行为而闻名于世（并因此在 1911 年被封为爵士）。他也是一位爱尔兰民族主义者。1914 至 1916 年间，凯斯门特秘密前往当时与英国交战的德国，试图招募爱尔兰战俘为爱尔兰独立而与英国作战。返回爱尔兰后，他被英国当局抓获，并以叛国罪受审。

审判凯斯门特所依据的法律是\textit{《1351 年叛国法》}。该法规定了何谓叛国罪，因此控方必须在审判中证明凯斯门特的行为符合《叛国法》所列的标准。相关段落规定，若某人

\begin{quote}
	在王国内，依附于国王的敌人，在王国内给予他们援助与慰藉，或在其他地方。
\end{quote}

则犯有叛国罪。

凯斯门特辩护的关键在于该句中的最后一个逗号，该逗号在 1351 年法律的原始法文文本中并不存在。凯斯门特“依附于国王的敌人”这一点并无争议，但问题在于，这种依附是否仅当发生在国内时才构成叛国，在国外发生是否也算叛国。辩方主张该法律条文存在歧义。其声称的歧义在于：“或在其他地方”是仅修饰“给予他们援助与慰藉”（无逗号时的自然解读），还是同时修饰“依附于国王的敌人”和“给予他们援助与慰藉”（有逗号时的自然解读）。尽管前一种解读可能显得牵强，但支持它的论证实际上并非没有说服力。然而，法庭最终裁定应按照有逗号的方式解读该段落，因此凯斯门特在德国的所作所为构成叛国，他被判处死刑。凯斯门特本人写道，他是“被一个逗号绞死的”。

我们可用真值函项逻辑对这一段落的不同解读进行符号化，借此消除歧义。首先，我们需要一套符号化约定：

\begin{ekey}
	\item[A] Casement 在国内支持国王的敌人
	\item[G] Casement 在国内给予国王的敌人援助和慰藉
	\item[B] Casement 在国外支持国王的敌人
	\item[H] Casement 在国外给予国王的敌人援助和慰藉
\end{ekey}

认为 Casement 行为不构成叛国的解读是：
\[A \lor (G \lor H)\]
而使他被定罪的解读则可以符号化为：
\[(A \lor B) \lor (G \lor H)\]
请记住，在我们讨论的这个案例中，Casement 在国外支持国王的敌人（$B$~为真），但在国内不支持，并且他无论国内国外均未给予国王的敌人援助或慰藉（$A$, $G$, 和~$H$ 均为假）。

英语中一种常见的结构歧义源于其缺少括号。例如，如果我说“我喜欢不又长又无聊的电影”，你很可能会认为我讨厌的是那些又长又无聊的电影。一种不太可能但逻辑上成立的解读是，我喜欢那些既（a）不长又（b）无聊的电影。第一种解读可能性更大，因为谁喜欢无聊的电影呢？但对于“我喜欢不甜而可口的菜肴”呢？在这里，更可能的解读是，我喜欢咸鲜可口的菜肴。（当然，我本可以表达得更清楚些，例如，“我喜欢不甜却可口的菜肴”。）类似的歧义也源于“且”与“或”的相互作用。例如，假设我请你给我发一张小而危险或隐秘的动物的照片。猎豹算吗？它很隐秘，但并不小。所以这取决于我是在找既小又危险或隐秘的动物（猎豹不算），还是在找要么小而危险的动物，要么是隐秘的动物（任意体型皆可）。

这类歧义被称为 \emph{辖域歧义}，因为它们取决于一个联结词是否在另一个联结词的辖域之内。例如，句子“《复仇者联盟：终局之战》不又长又无聊”，在以下两种解读之间是歧义的：
\begin{enumerate}
	\item\label{scamb1} 《复仇者联盟：终局之战》并非：既长又无聊。
	\item\label{scamb2} 《复仇者联盟：终局之战》既是：不长，又无聊。
\end{enumerate}

\Cref*{scamb2} 显然是假的，因为《复仇者联盟：终局之战》长达三个多小时。至于~\cref*{scamb1} 是否为真，则取决于这部电影是否无聊。我们可以使用以下符号化约定：
\begin{ekey}
	\item[B] 《复仇者联盟：终局之战》很无聊
	\item[L] 《复仇者联盟：终局之战》很长
\end{ekey}

现在，\Cref*{scamb1} 可以符号化为 “$\enot(L \eand B)$”，而 \cref*{scamb2} 则是 “$\enot L \eand B$”。在第一种情况下，“\eand”在“\enot”的辖域内；在第二种情况下，“\enot”在“\eand”的辖域内。

句子“太郎既小又危险或隐秘”在以下两种解读之间是歧义的：
\begin{enumerate}
	\item\label{scamb3} 太郎要么既小又危险，要么是隐秘的。
	\item\label{scamb4} 太郎既小，而且要么危险要么隐秘。
\end{enumerate}

我们可以使用以下符号化约定：
\begin{ekey}
	\item[D] 太郎是危险的
	\item[S] 太郎是小的
	\item[T] 太郎是隐秘的
\end{ekey}

\Cref*{scamb3} 的符号化是 “$(S \eand D) \eor T$”，而 \cref*{scamb4} 的符号化是 “$S \eand (D \eor T)$”。在第一种情况下，“\eand”在“\eor”的辖域内；在第二种情况下，“\eor”在“\eand”的辖域内。

\practiceproblems
\solutions
\problempart 以下句子是有歧义的。为每个句子提供符号化约定，并对不同解读进行符号化。
\begin{compactlist}
	\item Haskell 是一位观鸟者并且喜欢观看鹤。
	\item 动物园里有狮子或老虎和熊。
	\item 这花不红或香。
\end{compactlist}

\chapter{使用与提及}\label{s:UseMention}
在这部分，我们花了很多篇幅来\emph{讨论}语句。因此我们应该停下来解释一个重要且非常普遍的观点。

\section{引述约定}
考虑下面两个句子：
\begin{itemize}
		\item 亚里士多德是逻辑学的奠基人。
		\item 表达式“亚里士多德”由一个拉丁大写字母和八个拉丁小写字母组成。
	\end{itemize}

当我们想要谈论逻辑学的奠基人时，我们\emph{使用}他的名字。当我们想要谈论逻辑学奠基人的名字本身时，我们\emph{提及}那个名字，具体做法是将其放在引号内。

这里有一个一般性的要点。当我们想要谈论世界中的事物时，我们只是直接\emph{使用}词语。当我们想要谈论词语本身时，我们通常必须\emph{提及}这些词语。我们需要表明我们是在提及它们，而不是在使用它们。为此，我们需要一些约定。我们可以给它们加上引号，或者将它们居中显示（比如说）。所以下面这个句子：


\begin{itemize}
		\item “亚里士多德”是一位古希腊哲学家。
	\end{itemize}


说的是某个\emph{表达式}是一位古希腊哲学家。
这是假的。那个\emph{人}是一名希腊哲学家；他的\emph{名字}不是。相反，
下面这个句子：


\begin{itemize}
		\item 亚里士多德由一个大写字母和八个小写字母组成。
	\end{itemize}


也说了假的东西：亚里士多德是一个人（或者说在他活着的时候是），由血肉而非字母构成。最后一个例子：


\begin{itemize}
		\item “‘Aristotle’”是 Aristotle 的名称。
	\end{itemize}


这里，在左边，我们有一个名称的名称。在右边，我们有一个名称。也许这种句子只出现在逻辑教科书中，但它确实是真的。

这些只是一般的引用规则，你应该在所有作业中仔细遵守！需要明确的是，这里的引号不表示间接引语。它们表示你从谈论一个对象，转向谈论该对象的名称。

\section{对象语言与元语言}
这些一般的引用规则对我们来说非常重要。毕竟，我们在这里描述一种形式语言，即真值函项逻辑，所以我们必须经常\emph{提及}真值函项逻辑中的表达式。

当我们谈论一种语言时，我们正在谈论的那种语言被称为\define{对象语言}。我们用来谈论对象语言的语言被称为\define{元语言}。
\label{def.metalanguage}
\newglossaryentry{object language}
{
name=对象语言,
description={逻辑学家构造和研究的语言。在本教材中，对象语言是真值函项逻辑和一阶逻辑}
}

\newglossaryentry{metalanguage}
{
name=元语言,
description={逻辑学家用来谈论对象语言的语言。在本教材中，元语言是英语，辅以某些符号（如元变元）和术语（如“有效的”）}
}

在本章大部分情况下，对象语言一直是我们正在构建的形式语言：真值函项逻辑。元语言是英语。不完全是日常会话的英语，而是补充了一些额外词汇以方便我们行文的英语。

现在，我们使用大写字母作为真值函项逻辑的语句字母：
	$$A, B, C, Z, A_1, B_4, A_{25}, J_{375},\ldots$$
这些是对象语言（真值函项逻辑）的语句。它们不是英语句子。因此我们绝不能这样说，例如：


\begin{itemize}
		\item $D$ 是真值函项逻辑的一个语句字母。
	\end{itemize}


显然，我们试图用英语句子来谈论对象语言（真值函项逻辑），但“$D$”是真值函项逻辑的一个语句，而不是英语的一部分。因此，上述说法是胡言乱语，就像：


\begin{itemize}
		\item Schnee ist weiß 是一个德语句子。
	\end{itemize}


在这种情况下，我们真正想说的是：


\begin{itemize}
		\item “Schnee ist weiß”是一个德语句子。
	\end{itemize}


同样，我们上面真正想说的只是：


\begin{itemize}
		\item “$D$”是真值函项逻辑的一个语句字母。
	\end{itemize}


一般要点是，每当我们想用英语谈论真值函项逻辑中的某个特定表达式时，我们需要表明我们是在\emph{提及}这个表达式，而不是\emph{使用}它。我们可以使用引号，或者采用一些类似的约定，例如将其放在页面中央。

\section{元变元}\label{s:Metavariables}
然而，我们不仅仅想谈论真值函项逻辑的\emph{特定}表达式。我们还希望能够谈论真值函项逻辑的\emph{任意}语句。事实上，我们在\cref{s:Sentences}中已经这样做了，当时我们给出了真值函项逻辑语句的归纳定义。我们使用大写花体字母来做到这一点，即：
	$$\metav{A}, \metav{B}, \metav{C}, \metav{D}, \ldots$$
这些符号本身并非真值函项逻辑的一部分。相反，它们是我们用来谈论真值函项逻辑\emph{任意}表达式的（扩充后的）元语言的一部分。为了解释为什么我们需要它们，回忆一下真值函项逻辑语句归纳定义的第二条：


\begin{compactlist}
		\item[2.] 如果 $\metav{A}$ 是一个语句，那么 $\enot \metav{A}$ 是一个语句。
	\end{compactlist}


这里谈论的是\emph{任意}语句。如果我们改为：


\begin{compactlist}
		\item[2$'$.] 如果 “$A$” 是一个语句，那么 “$\enot A$” 是一个语句。
	\end{compactlist}

这就没法让我们判断“$\enot B$”是不是一个语句。需要强调的是：
	\factoidbox{
	  “$\metav{A}$”是扩充英语中的一个符号（称为\define{元变元（metavariable）}），我们用它来谈论真值函项逻辑的表达式。	“$A$”是真值函项逻辑中的一个特定的语句字母。}

        \newglossaryentry{metavariables}
{
name=元变元,
description={元语言中的一个变元，可以代表对象语言中的任意语句}
}
但这最后一个例子引出了一个进一步的复杂性，涉及引用约定。在我们归纳定义的第二条中，我们完全没有使用任何引号。我们本应该用吗？

问题在于，这条规则右边的表达式，即“$\enot\metav{A}$”，并不是一个英语句子，因为它包含了“$\enot$”。所以我们或许可以尝试写：


\begin{numberlist}
		\item[2$''$.] 如果 \metav{A} 是一个语句，那么“$\enot \metav{A}$”就是一个语句。
	\end{numberlist}


但这并不正确：“$\enot \metav{A}$”并不是真值函项逻辑语句，因为“$\metav{A}$”是（扩充）英语的符号，而不是真值函项逻辑的符号。

我们真正想说的是：


\begin{numberlist}
		\item[2$'''$.] 如果 \metav{A} 是一个语句，那么把符号“$\enot$”和语句 \metav{A} 拼接起来的结果就是一个语句。
	\end{numberlist}


这完全正确，但相当冗长。%Quine introduced a convention that speeds things up here. In place of (2$''$), he suggested:
%	\begin{numberlist}
%		\item[2$'''$.] If \metav{A} and \metav{B} are sentences, then $\ulcorner (\metav{A}\eand\metav{B})\urcorner$ is a sentence
%	\end{numberlist}
%The rectangular quote-marks are sometimes called “Quine quotes', after Quine. The general interpretation of an expression like ”$\ulcorner (\metav{A}\eand\metav{B})\urcorner$” is in terms of rules for concatenation.
但我们可以通过建立自己的约定来避免冗长。我们完全可以规定，像“$\enot \metav{A}$”这样的表达式就应该\emph{直接}依据拼接规则来理解。所以，\emph{正式说来}，元语言表达式“$\enot \metav{A}$”
是下面这个说法的简写：


\begin{quote}
	把符号“$\enot$”和语句 \metav{A} 拼接起来的结果
\end{quote}


类似地，对于“$(\metav{A} \eand \metav{B})$”、“$(\metav{A} \eor \metav{B})$”等表达式也是如此。

\section{论证的引用约定}
我们使用真值函项逻辑的主要目的之一就是研究论证，这将是\cref{ch.TruthTables}中要讨论的内容。在英语中，一个论证的前提通常由若干个独立的句子表达，结论则由另一个句子表达。既然我们可以符号化英语句子，我们也就能用真值函项逻辑来符号化英语论证。

或者更准确地说，我们可以用真值函项逻辑来符号化一个英语论证中使用的每一个\emph{语句}。然而，真值函项逻辑本身没有办法标明其中哪些语句是\emph{前提}，哪个是论证的\emph{结论}。（这与自然英语形成对比，英语会用“所以”、“因此”等词语来标记一个句子是论证的\emph{结论}。）

%So, if we want to symbolize an \emph{argument} in TFL, what are we to do?

%An obvious thought would be to add a new symbol to the \emph{object} language of TFL itself, which we could use to separate the premises from the conclusion of an argument. However, adding a new symbol to our object language would add significant complexity to that language, since that symbol would require an official syntax.\footnote{\emph{The following footnote should be read only after you have finished the entire book!} And it would require a semantics. Here, there are deep barriers concerning the semantics. First: an object-language symbol which adequately expressed “therefore' for TFL would not be truth-functional. (\emph{Exercise}: why?) Second: a paradox known as ”validity Curry' shows that FOL itself \emph{cannot} be augmented with an adequate, object-language “therefore'.}

因此，我们需要再引入一点记号。假设我们要用 $\metav{A}_1$、……、$\metav{A}_n$ 来符号化一个论证的前提，用 $\metav{C}$ 来符号化其结论。那么我们就写作：
$$\metav{A}_1, \ldots, \metav{A}_n \therefore \metav{C}$$
符号“$\therefore$”的作用仅仅是标明哪些语句是前提，哪个是结论。

%Strictly, this extra notation is \emph{unnecessary}. After all, we could always just write things down long-hand, saying: the premises of the argument are well symbolized by $\metav{A}_1, \ldots \metav{A}_n$, and the conclusion of the argument is well symbolized by $\metav{C}$. But having some convention will save us some time. Equally, the particular convention we chose was fairly \emph{arbitrary}. After all, an equally good convention would have been to underline the conclusion of the argument. Still, this is the convention we will use.

严格说来，符号“$\therefore$”并不属于对象语言，而是属于\emph{元语言}。因此，有人可能会认为，我们需要在它两侧的真值函项逻辑语句周围加上引号。这个想法是合理的，但加上这些引号会使文本更难阅读。而且——与上文一样——别忘了是\emph{我们}在规定一些新约定。所以，我们完全可以规定这些引号是不必要的。也就是说，我们可以直接写作：
$$A, A \eif B \therefore B$$
\emph{不加任何引号}，来表示一个其前提是（由）“$A$”和“$A \eif B$”符号化、结论是（由）“$B$”符号化的论证。
